\section{更广泛的应用场景}

AutoResearch 的适用范围远不止 Claude skill。作者明确指出：\textbf{任何你能打分的东西，都能用这套方法。}

\subsection{具体案例}

\begin{itemize}
    \item \textbf{网站速度优化}：改一处代码 $\to$ 测量页面加载时间 $\to$ 留下或撤销。有人在 67 轮迭代后将页面加载时间\textbf{从 1100ms 降到 67ms}。

    \item \textbf{Cold outreach（陌生客户开发邮件）}：定义 checklist——``有没有提到对方公司？是否在 75 字以内？是否以具体问题收尾？''让 agent 跑 50 个变体。

    \item \textbf{Newsletter 开篇}：``开场白有没有包含个人细节？'' ``有没有使用陈词滥调？''让 agent 在 autopilot 上打磨你的文字。

    \item \textbf{任何你反复使用的 prompt}——只要能定义打分标准，就能跑 autoresearch。
\end{itemize}

\begin{importantbox}{通用原则}
\textbf{If you can score it, you can autoresearch it.}

关键前提只有一个：你的优化目标必须是可量化的。能用是/否 checklist 打分，就能用 autoresearch 自动迭代。
\end{importantbox}

\subsection{本章小结}

AutoResearch 方法的唯一前提是``可量化''。从网站性能到邮件写作、从 newsletter 到任何反复使用的 prompt，只要能定义出是/否评分标准，就能套用这套自动迭代框架。

